\chapter[Imitation learning and offline RL]{Imitation learning and offline
  reinforcement learning}
\label{ch:imitation}

The previous chapters assumed that the agent learns by interacting with the
environment and receiving rewards. In many applications this is expensive
or dangerous (a car, a patient, a power grid), and the reward is hard to
write down, while \emph{data} are available: demonstrations by an expert,
or logs of past decisions. This chapter studies two ways of using them.
\emph{Imitation learning} learns from demonstrations, without rewards;
\emph{offline reinforcement learning} learns from a fixed dataset of
transitions with rewards, without further interaction.

\section{Behaviour cloning}

The simplest approach treats demonstrations $\mathcal{D} = \{(s_i, a_i)\}$
of an expert policy $\pi^\star$ as a supervised learning problem: fit
$\pi_\theta$ by maximum likelihood,
\begin{equation}
  \min_\theta\; \E_{(s, a) \sim \mathcal{D}}\left[ -\log \pi_\theta(a \mid
  s) \right].
\end{equation}
Behaviour cloning (BC) is simple and often works well with enough data
(it was used to drive a car as early as 1989, by ALVINN). Its weakness is
that the states of the training set are those visited by the expert. When
the learner makes a mistake, it reaches states the expert never visited,
where it has no data and makes more mistakes: the errors compound.

\begin{theorem}{Compounding errors of behaviour cloning}{bc}
  Consider episodes of $T$ steps with costs $c(s, a) \in [0, 1]$, a
  deterministic expert $\pi^\star$, and let $J(\pi)$ be the expected total
  cost of $\pi$. Let $d_t$ be the distribution of the expert's state at step
  $t$, and suppose that the learner $\pi$ disagrees with the expert with
  probability at most $\varepsilon$ on these distributions:
  $\Prob_{s \sim d_t,\, a \sim \pi(\cdot \mid s)}(a \neq \pi^\star(s)) \leq
  \varepsilon$ for all $t$. Then
  \begin{equation}
    J(\pi) \leq J(\pi^\star) + \varepsilon\, T^2.
  \end{equation}
\end{theorem}

\begin{proof}
  Run the expert and the learner with the same random numbers for the
  environment, so that their trajectories coincide as long as their actions
  coincide. Let $D_t$ be the event that the trajectories have diverged at
  or before step $t$. If they have not diverged before step $t$, the shared
  state is the expert's state, distributed as $d_t$, so the probability of
  a first disagreement at step $t$ is at most $\varepsilon$, and $\Prob(D_t)
  \leq t \varepsilon$. At step $t$ the costs of the two trajectories are
  equal outside $D_t$ and differ by at most $1$ on $D_t$. Summing,
  $J(\pi) - J(\pi^\star) \leq \sum_{t=1}^T t\varepsilon \leq \varepsilon
  T^2$.
\end{proof}

\begin{intuition}
  A learner who has only watched an expert drive has never seen the car
  on the shoulder of the road. After a small mistake it is in a situation
  it never trained on, where it makes more mistakes, and it may never
  return to familiar ground. A mistake at step $t$ can spoil all the
  remaining steps, so the cost of a per-step error rate $\varepsilon$ is
  multiplied by the horizon \emph{twice}.
\end{intuition}

The bound is tight up to constants: there are problems where a single
mistake leads to states from which the learner never recovers, and the
cost is then of order $\varepsilon T^2$, not the $\varepsilon T$ of a
supervised learning error accumulated over $T$ steps.

\section{DAgger}

The cure is to train the learner on the states \emph{it} visits. DAgger
(dataset aggregation, Ross, Gordon and Bagnell, 2011) requires an expert
that can be queried in any state.

\begin{algobox}{DAgger}
  \begin{algorithmic}[1]
    \State $\mathcal{D} \gets$ demonstrations of the expert; train
      $\pi_1$ on $\mathcal{D}$
    \For{$i = 1, 2, \dots, N$}
      \State run $\pi_i$ (optionally mixed with the expert) to collect states
        $s$
      \State ask the expert for its actions $\pi^\star(s)$; $\mathcal{D}
        \gets \mathcal{D} \cup \{(s, \pi^\star(s))\}$
      \State train $\pi_{i+1}$ on $\mathcal{D}$
    \EndFor
  \end{algorithmic}
\end{algobox}

DAgger is an online learning algorithm in disguise: training on the
aggregated data is ``follow the leader'' on the sequence of losses defined
by the learner's own state distributions. Ross et al.\ show, with a
no-regret learner, that some policy of the sequence has a cost $J(\pi) \leq
J(\pi^\star) + u\, T \varepsilon_N + o(1)$, where $\varepsilon_N$ is the
best training error on the aggregated data and $u$ bounds the increase of
the expert's cost-to-go after one wrong action: the error grows linearly in
$T$ when the expert can recover from the learner's mistakes.

\begin{intuition}
  DAgger is a driving instructor in the passenger seat: the learner
  drives, and in the situations \emph{it} gets into, including those
  caused by its own mistakes, the instructor says what it should have
  done. The training data then cover the states the learner actually
  visits, and the learner learns how to recover.
\end{intuition}

\begin{pacman}{Expert iteration is DAgger}{dagger}
  In \cref{ch:model-based}, the expert is the look-ahead search and the
  learner the PPO network: \code{pacman\_expert.py} lets the network play
  and asks the search for its move in the states the network reached, which
  is DAgger. The search can be queried in any state because it only needs
  a copy of the game. Training on games played \emph{by the search} would
  have been behaviour cloning, exposed to \cref{thm:bc}.
\end{pacman}

\section{Inverse reinforcement learning}

Behaviour cloning copies actions; inverse reinforcement learning (IRL)
recovers the \emph{reward} that the expert optimises, which transfers to new
dynamics and explains the behaviour. The problem is ill-posed: the reward
$r = 0$ makes every policy optimal, and adding to the reward a
\emph{shaping} term $\gamma \Phi(s') - \Phi(s)$, for any function $\Phi$,
changes the reward without changing the optimal policies (Ng, Harada and
Russell, 1999).
A principle is needed to choose among the rewards that explain the data.

\paragraph{Maximum entropy IRL.} Ziebart et al.\ (2008) take a linear
reward $r_\theta(\tau) = \theta^\top \phi(\tau)$ of the features $\phi(\tau)
= \sum_t \phi(s_t, a_t)$ of a trajectory, and model the expert as choosing
trajectories with probability $p_\theta(\tau) = \exp(\theta^\top
\phi(\tau)) / Z(\theta)$ (for deterministic dynamics): the maximum entropy
distribution that matches the feature expectations. The parameters are
fitted by maximum likelihood on the demonstrations.

\begin{proposition}{Gradient of maximum entropy IRL}{maxent-irl}
  The average log-likelihood $L(\theta) = \frac{1}{N} \sum_{i=1}^N \log
  p_\theta(\tau_i)$ of $N$ demonstrations is concave, with gradient
  \begin{equation}
    \nabla L(\theta) = \frac{1}{N} \sum_{i=1}^N \phi(\tau_i) -
    \E_{\tau \sim p_\theta}\left[ \phi(\tau) \right].
  \end{equation}
\end{proposition}

\begin{proof}
  $L(\theta) = \theta^\top \bar\phi - \log Z(\theta)$ with $\bar\phi$ the
  empirical mean of the features, and $\nabla \log Z(\theta) =
  \sum_\tau \phi(\tau) e^{\theta^\top \phi(\tau)} / Z(\theta) =
  \E_{p_\theta}[\phi]$. The Hessian of $\log Z$ is the covariance of $\phi$
  under $p_\theta$, positive semi-definite, so $L$ is concave.
\end{proof}

\begin{intuition}
  The gradient compares what the expert does with what the current model
  of the expert does. A feature the expert's trajectories show more often
  (staying far from the ghosts, say) gets a higher reward weight; a feature
  the model overproduces gets a lower one. Learning stops when the model's
  trajectories have the same feature statistics as the expert's.
\end{intuition}

At the optimum the model matches the expert's feature expectations. The
expectation under $p_\theta$ requires solving a soft (maximum entropy) RL
problem for the current reward, as in SAC (\cref{ch:continuous}): IRL is
an outer loop around RL, and therefore expensive.

\paragraph{Adversarial imitation.} GAIL (Ho and Ermon, 2016) skips the
reward: a discriminator $D_\psi(s, a)$ is trained to distinguish the
expert's pairs from the learner's, and the policy is trained by RL (TRPO or
PPO) with the reward $-\log(1 - D_\psi(s, a))$, so that it produces pairs
the discriminator takes for the expert's. As in generative adversarial
networks, at the equilibrium the occupancy measure of the learner matches
that of the expert (the procedure minimises their Jensen--Shannon
divergence). Unlike BC, GAIL matches distributions of states, not only
actions in the expert's states, and does not suffer from compounding
errors; but it needs interaction with the environment.

\section{Offline reinforcement learning}

In offline (or batch) RL, the agent receives a dataset $\mathcal{D} =
\{(s_i, a_i, r_i, s'_i)\}$ collected by one or several \emph{behaviour}
policies $\beta$, and must produce the best possible policy without
collecting more data. Q-learning and DQN are off-policy, so it is tempting
to run them on $\mathcal{D}$. This fails in practice, often badly.

\paragraph{Extrapolation error.} The target $r + \gamma \max_{a'} Q(s',
a')$ evaluates $Q$ at actions $a'$ that may never appear in the data for
$s'$. There the network's values are arbitrary, and the maximisation
selects precisely the actions whose values are overestimated. Online, the
agent would try these actions and correct the error; offline, nothing
corrects it, and the errors propagate through the targets. The
\emph{distribution shift} between the learnt policy and $\beta$ is the
central problem of offline RL. Three families of methods address it.

\paragraph{Policy constraints.} Keep the policy close to the behaviour
policy. BCQ (Fujimoto et al., 2019) only maximises over actions proposed by
a generative model of $\beta$. TD3+BC (Fujimoto and Gu, 2021) adds a
behaviour cloning term to the actor of TD3 (\cref{ch:continuous}):
\begin{equation}
  \max_\theta\; \E_{(s, a) \sim \mathcal{D}}\left[ \lambda\, Q_\phi(s,
  \mu_\theta(s)) - \norm{\mu_\theta(s) - a}^2 \right],
\end{equation}
a one-line change that is a strong baseline.

\paragraph{Conservative values.} CQL (Kumar et al., 2020) learns a
Q-function that underestimates the values of the actions absent from the
data. For discrete actions it adds to the TD loss the term
\begin{equation}
  \label{eq:cql}
  \alpha\, \E_{s \sim \mathcal{D}}\left[ \log \sum_{a} \exp Q(s, a) -
  \E_{a \sim \mathcal{D}(\cdot \mid s)}[Q(s, a)] \right],
\end{equation}
which pushes down the values of all actions (through a soft maximum) and
pushes up those of the actions in the data.

\begin{lstlisting}[caption={CQL loss for discrete actions (PyTorch)}]
q = q_net(states)                                  # (B, |A|)
q_data = q.gather(1, actions[:, None]).squeeze(1)  # Q(s, a), a in data
with torch.no_grad():
    target = rewards + gamma * (1 - dones) * \
        target_net(next_states).max(1).values
td_loss = F.smooth_l1_loss(q_data, target)
cql_loss = (torch.logsumexp(q, dim=1) - q_data).mean()
loss = td_loss + alpha * cql_loss
\end{lstlisting}

The following idealised, tabular version shows why the values become
conservative. Let $\mu(\cdot \mid s)$ be the policy whose values are pushed
down (in \eqref{eq:cql}, the softmax of $Q$), and consider the iteration
\begin{equation*}
  \hat Q_{k+1} = \argmin_Q\; \alpha \left( \E_{s, a \sim \mu} [Q(s, a)] -
  \E_{s, a \sim \beta}[Q(s, a)] \right) + \frac{1}{2} \E_{s, a \sim \beta}
  \left[ \left( Q(s, a) - (\T^\pi \hat Q_k)(s, a) \right)^2 \right],
\end{equation*}
where the states are drawn from the data and the pairs weighted by
$\beta(a \mid s) > 0$.

\begin{proposition}{CQL underestimates the value}{cql}
  The minimiser is $\hat Q_{k+1}(s, a) = (\T^\pi \hat Q_k)(s, a) - \alpha
  \left( \frac{\mu(a \mid s)}{\beta(a \mid s)} - 1 \right)$, and
  \begin{equation}
    \E_{a \sim \mu(\cdot \mid s)}\left[ \hat Q_{k+1}(s, a) \right] \leq
    \E_{a \sim \mu(\cdot \mid s)}\left[ (\T^\pi \hat Q_k)(s, a) \right].
  \end{equation}
\end{proposition}

\begin{proof}
  The objective separates over the pairs $(s, a)$. For one pair, the
  derivative with respect to $Q(s, a)$ is $\alpha (\mu(a \mid s) - \beta(a
  \mid s)) + \beta(a \mid s) (Q(s, a) - (\T^\pi \hat Q_k)(s, a))$ (up to
  the common factor of the state's weight); setting it to zero gives the
  minimiser. For the inequality, $\E_\mu[\mu / \beta - 1] = \sum_a \mu(a)^2
  / \beta(a) - 1 \geq 0$, by the Cauchy--Schwarz inequality $1 = (\sum_a
  \mu(a))^2 \leq \sum_a \mu(a)^2 / \beta(a) \cdot \sum_a \beta(a)$.
\end{proof}

The penalty is large where $\mu$ puts much more weight than the data:
exactly the actions whose values cannot be trusted.

\begin{intuition}
  CQL tells the Q-function: ``be pessimistic about what you have not
  seen''. An action that the policy would choose much more often than the
  data contain it gets its value pushed down, in proportion to that
  imbalance; actions well covered by the data are left almost unchanged.
  The policy then prefers actions whose good value is backed by evidence.
\end{intuition}

\paragraph{In-sample learning.} IQL (Kostrikov et al., 2022) never
evaluates $Q$ at actions outside the data. It fits $V(s)$ to an upper
expectile of $Q(s, a)$ over the actions of the data, minimising
$\E_{(s, a) \sim \mathcal{D}}[L_\tau(Q(s, a) - V(s))]$ with $L_\tau(u) =
\abs{\tau - \mathbb{1}(u < 0)}\, u^2$ and $\tau \approx 0.7$--$0.9$, which
approximates a maximum over the actions supported by the data; $Q$ is
fitted to $r + \gamma V(s')$; and the policy is extracted by
advantage-weighted regression, maximising $\E_{\mathcal{D}}[e^{b (Q(s, a) -
V(s))} \log \pi(a \mid s)]$: behaviour cloning weighted by the advantage.

\begin{intuition}
  An upper expectile of $Q(s, a)$ over the actions of the data is a soft
  version of ``the best action that was actually tried in this state''. IQL
  thus computes a maximum over the data's actions only, never asking the
  network about actions it has no data for, and then imitates the data,
  copying the good actions more than the bad ones.
\end{intuition}

\begin{remark}{Offline RL and Pacman}{offline-pacman}
  The replay memory of a DQN training run is an offline dataset of
  millions of transitions of Pacman. Running DQN on it from scratch, without
  interaction, is an instructive experiment: the values of the rare
  actions (walking into a ghost) are poorly constrained by the data, and
  the conservative loss above is a few lines to add to \code{DQN.learn}.
\end{remark}

\section{Summary}

\begin{itemize}[nosep]
  \item Behaviour cloning is supervised learning on demonstrations; its
    errors compound quadratically with the horizon, because the learner
    leaves the states covered by the data.
  \item DAgger queries the expert in the learner's own states and reduces
    the error to linear in the horizon; expert iteration with a search as
    expert is DAgger.
  \item Inverse RL recovers a reward (maximum entropy IRL matches feature
    expectations); GAIL matches occupancy measures adversarially.
  \item Offline RL suffers from extrapolation error at actions absent from
    the data; policy constraints (TD3+BC), conservative values (CQL) and
    in-sample learning (IQL) avoid it.
\end{itemize}
